\documentclass[11pt]{article}

\usepackage[margin=1in]{geometry}
\usepackage{amsmath, amssymb, amsthm}
\usepackage{enumitem}

\pagestyle{empty}

\begin{document}

\begin{center}
  {\Large\bfseries Weekly Assignment: Formal Languages \&\\[2pt]
   Regular Expressions (Epp~12.1)}
\end{center}

\bigskip

\noindent
Complete each of the following problems.  Show all work and justify your answers.
Recall the order of precedence for regular expression operations:
$*$ (highest), then concatenation, then $|$ (lowest).
Also recall the shorthand notations: character classes such as
$[0\text{--}9]$, the $+$ operator ($r^+ = rr^*$), the $?$
operator ($r? = (\lambda | r)$), and $r\{n\}$ for exactly $n$
repetitions.

\bigskip

%% ─────────────────────────────────────────────
\noindent\textbf{Problem 1} (Epp 12.1 \#2). \;
Let $\Sigma = \{x, y\}$ be an alphabet.

\begin{enumerate}[label=\textbf{\alph*.},itemsep=4pt]
  \item Let $L_3$ be the language consisting of all strings
        over~$\Sigma$ of length $\le 3$ in which all the
        $x$'s appear to the left of all the $y$'s.  List the
        elements of~$L_3$ between braces.
  \item List between braces the elements of~$\Sigma^4$, the
        set of all strings of length~4 over~$\Sigma$.
  \item Let $A = \Sigma^0 \cup \Sigma^1$ and
        $B = \Sigma^3 \cup \Sigma^4$.  Describe $A$, $B$, and
        $A \cup B$ in words.
\end{enumerate}

\bigskip

%% ─────────────────────────────────────────────
\noindent\textbf{Problem 2} (Epp 12.1 \#11). \;
Use the rules about order of precedence to eliminate the
parentheses in the following regular expression:
\[
  \bigl(1(1^*)\bigr) \;|\;
  \Bigl(\bigl(1(0^*)\bigr) \;|\; \bigl((1^*)1\bigr)\Bigr)
\]

\bigskip

%% ─────────────────────────────────────────────
\noindent\textbf{Problem 3} (Epp 12.1 \#23). \;
Indicate whether the given strings belong to the language
defined by the regular expression below.  Briefly justify
your answers.
\[
  \text{Expression: } (x^*y | zy^*)^*
  \qquad\qquad
  \text{Strings: } zyyxz,\;\; zyyzy
\]

\bigskip

%% ─────────────────────────────────────────────
\noindent\textbf{Problem 4} (Epp 12.1 \#38). \;
Write a regular expression to define the set of all telephone
numbers with the following format: three digits, then a
hyphen, then three more digits, then a hyphen, and then four
digits, where the first three digits are either \texttt{800}
or~\texttt{888} and the last four digits start and end
with a~\texttt{2}.

\end{document}
